% GENERATED FILE. Edit canonical lessons, then run npm run generate:books.
% canonical-source-hash: fnv1a64:d77dc65af6208ee1
% canonical-lessons: PT-C447-marginal, PT-C447-vaguear, PT-C447-espreitar, PT-C447-cirandar, PT-C447-miradouro

\chapter{Dar uma volta}
\label{ch:pt-dar-uma-volta}
\hlchaptermodality{447}

\begin{chapteropening}
\textbf{By the end of this chapter:} \emph{I can talk about strolling round a town: the seafront, roaming, peeking, pottering about and the viewpoint.}

Complete the last lesson of chapter 447: I can talk about strolling round a town: the seafront, roaming, peeking, pottering about and the viewpoint.
\end{chapteropening}

\section[a marginal]{a marginal — the seafront road, the riverside road}

\label{lesson:PT-C447-marginal}



\begin{quote}
Before the new one: say the Portuguese for the patron saint, then the Portuguese for to parade.
\end{quote}

\subsection*{You'll want to know: a marginal}
\textbf{a marginal} — \textquotedblleft{}the seafront road, the riverside road\textquotedblright{}.

\begin{quote}
Ao domingo, passeamos pela marginal.
\end{quote}

\begin{quote}
\emph{On Sundays, we walk along the seafront.}
\end{quote}

Also an adjective: \emph{marginal}, on the edge.

\begin{cousinweb}
\emph{marginal} is a learned borrowing from Medieval Latin \emph{marginālis} (EtymDB), from \emph{margō} 'edge, border': the road along the water's edge.
\end{cousinweb}

\begin{grammarlens}[title={What you've built}]
One more word for telling, explaining and describing at B1.
\end{grammarlens}

\subsection*{Your turn}
\begin{itemize}
\raggedright
  \item \emph{Say it:} \emph{a marginal}
  \item \emph{Say it:} \emph{a marginal}, once more
  \item \emph{Say it:} a sentence of your own with \emph{a marginal}
  \item \emph{Recall:} say the Portuguese for the patron saint, then the Portuguese for to parade, then say \emph{a marginal} again
\end{itemize}

\begin{tcolorbox}[breakable,colback=teal!4,colframe=teal!35!black,title={Before you move on}]
What does a marginal mean? (\textquotedblleft{}The seafront road, the riverside road\textquotedblright{}.) Say it once more.
\end{tcolorbox}
\section[vaguear]{vaguear — to roam, to wander}

\label{lesson:PT-C447-vaguear}



\begin{quote}
Before the new one: say the Portuguese for to parade, then the Portuguese for the seafront road.
\end{quote}

\subsection*{You'll want to know: vaguear}
\textbf{vaguear} — \textquotedblleft{}to roam, to wander\textquotedblright{}.

\begin{quote}
Vagueámos pela cidade sem destino.
\end{quote}

\begin{quote}
\emph{We roamed the town with no destination.}
\end{quote}

Thoughts can \emph{vaguear} too.

\begin{cousinweb}
\emph{vaguear} is built on \emph{vago} 'wandering, vague', which is a learned borrowing from Latin \emph{vagus} 'wandering' (EtymDB).
\end{cousinweb}

\begin{grammarlens}[title={What you've built}]
One more word for telling, explaining and describing at B1.
\end{grammarlens}

\subsection*{Your turn}
\begin{itemize}
\raggedright
  \item \emph{Say it:} \emph{vaguear}
  \item \emph{Say it:} \emph{vaguear}, once more
  \item \emph{Say it:} a sentence of your own with \emph{vaguear}
  \item \emph{Recall:} say the Portuguese for to parade, then the Portuguese for the seafront road, then say \emph{vaguear} again
\end{itemize}

\begin{tcolorbox}[breakable,colback=teal!4,colframe=teal!35!black,title={Before you move on}]
What does vaguear mean? (\textquotedblleft{}To roam, to wander\textquotedblright{}.) Say it once more.
\end{tcolorbox}
\section[espreitar]{espreitar — to peek, to peep}

\label{lesson:PT-C447-espreitar}



\begin{quote}
Before the new one: say the Portuguese for the seafront road, then the Portuguese for to roam.
\end{quote}

\subsection*{You'll want to know: espreitar}
\textbf{espreitar} — \textquotedblleft{}to peek, to peep\textquotedblright{}.

\begin{quote}
As crianças espreitam pela janela.
\end{quote}

\begin{quote}
\emph{The children peek through the window.}
\end{quote}

\emph{Espreitar uma oportunidade} — to watch for a chance.

\begin{cousinweb}
\emph{espreitar} is inherited from Old Portuguese \emph{espreitar} (EtymDB); it is usually traced to Latin \emph{explicitāre} 'to unfold, to examine', but that is not settled.
\end{cousinweb}

\begin{grammarlens}[title={What you've built}]
One more word for telling, explaining and describing at B1.
\end{grammarlens}

\subsection*{Your turn}
\begin{itemize}
\raggedright
  \item \emph{Say it:} \emph{espreitar}
  \item \emph{Say it:} \emph{espreitar}, once more
  \item \emph{Say it:} a sentence of your own with \emph{espreitar}
  \item \emph{Recall:} say the Portuguese for the seafront road, then the Portuguese for to roam, then say \emph{espreitar} again
\end{itemize}

\begin{tcolorbox}[breakable,colback=teal!4,colframe=teal!35!black,title={Before you move on}]
What does espreitar mean? (\textquotedblleft{}To peek, to peep\textquotedblright{}.) Say it once more.
\end{tcolorbox}
\section[cirandar]{cirandar — to potter about, to go here and there}

\label{lesson:PT-C447-cirandar}



\begin{quote}
Before the new one: say the Portuguese for to roam, then the Portuguese for to peek.
\end{quote}

\subsection*{You'll want to know: cirandar}
\textbf{cirandar} — \textquotedblleft{}to potter about, to go here and there\textquotedblright{}.

\begin{quote}
Passei a manhã a cirandar pelas lojas.
\end{quote}

\begin{quote}
\emph{I spent the morning pottering round the shops.}
\end{quote}

Built on \emph{a ciranda}, a sieve and also a round dance.

\begin{cousinweb}
\emph{cirandar} is built on \emph{ciranda} 'sieve; a round dance' (EtymDB lists the word); the origin of \emph{ciranda} is not settled.
\end{cousinweb}

\begin{grammarlens}[title={What you've built}]
One more word for telling, explaining and describing at B1.
\end{grammarlens}

\subsection*{Your turn}
\begin{itemize}
\raggedright
  \item \emph{Say it:} \emph{cirandar}
  \item \emph{Say it:} \emph{cirandar}, once more
  \item \emph{Say it:} a sentence of your own with \emph{cirandar}
  \item \emph{Recall:} say the Portuguese for to roam, then the Portuguese for to peek, then say \emph{cirandar} again
\end{itemize}

\begin{tcolorbox}[breakable,colback=teal!4,colframe=teal!35!black,title={Before you move on}]
What does cirandar mean? (\textquotedblleft{}To potter about, to go here and there\textquotedblright{}.) Say it once more.
\end{tcolorbox}
\section[o miradouro]{o miradouro — the viewpoint, the lookout}

\label{lesson:PT-C447-miradouro}



\begin{quote}
Before the new one: say the Portuguese for to peek, then the Portuguese for to potter about.
\end{quote}

\subsection*{You'll want to know: o miradouro}
\textbf{o miradouro} — \textquotedblleft{}the viewpoint, the lookout\textquotedblright{}.

\begin{quote}
Do miradouro vemos o rio e as pontes.
\end{quote}

\begin{quote}
\emph{From the viewpoint we can see both the river and the bridges.}
\end{quote}

In Brazil: \emph{o mirante}.

\begin{cousinweb}
\emph{miradouro} is built on \emph{mirar} 'to gaze', which is inherited from Vulgar Latin \emph{mīrāre} (EtymDB), from Latin \emph{mīrārī} 'to wonder at'.
\end{cousinweb}

\begin{grammarlens}[title={What you've built}]
One more word for telling, explaining and describing at B1.
\end{grammarlens}

\subsection*{Your turn}
\begin{itemize}
\raggedright
  \item \emph{Say it:} \emph{o miradouro}
  \item \emph{Say it:} \emph{o miradouro}, once more
  \item \emph{Say it:} a sentence of your own with \emph{o miradouro}
  \item \emph{Recall:} say the Portuguese for to peek, then the Portuguese for to potter about, then say \emph{o miradouro} again
\end{itemize}

\begin{tcolorbox}[breakable,colback=teal!4,colframe=teal!35!black,title={Before you move on}]
What does o miradouro mean? (\textquotedblleft{}The viewpoint, the lookout\textquotedblright{}.) Say it once more.
\end{tcolorbox}
